\section*{Synthèse}
\addcontentsline{toc}{section}{Synthèse}

rockim est un code de calcul écrit pendant la thèse, en C++ avec la seule bibliothèque Eigen, pour simuler l'impact d'un insert en carbure sur une roche. Il implémente la méthode combinée éléments finis et éléments discrets (FDEM) de la lignée de Munjiza~\cite{munjiza2004fdem} : la roche est découpée en tétraèdres élastiques reliés par des joints cohésifs qui s'endommagent puis rompent, et les fragments libérés interagissent par contact frottant. Depuis la décision du 14 août 2026, le développement porte sur les modes FDEM 2D et 3D ; les modes éléments finis continus et éléments discrets sont gelés, et la production 3D calibrée de la thèse reste confiée à Abaqus avec les VUMAT de la thèse.

Le code résout correctement plusieurs de ses briques, et ces vérifications se rejouent en une commande. Le contact par potentiel conserve l'énergie d'une collision élastique à $2\times10^{-8}$ près en 3D ; l'algèbre du contact d'outil en vitesse est juste à $2{,}2\times10^{-16}$ ; une barre en traction donne l'allongement $\sigma L/E$ à $0{,}03~\%$ ; les 26 contrôles à charge nulle ne rompent aucun joint ; le bilan d'énergie à sept postes ferme à l'arrondi en élastique. Une onde de compression dans une barre est reproduite par le solveur éléments finis 3D avec une célérité juste à $0{,}05~\%$ et un ordre de convergence observé de $1{,}58$. Un bloc qui glisse jusqu'à l'arrêt s'arrête à $1{,}74~\%$ de la distance exacte avec le contact par potentiel.

Ces mêmes bancs ont révélé deux biais qui touchent directement l'impact. Avec des joints présents dès l'origine et une pénalité de 20 fois $E/h$, l'onde est ralentie de $2{,}9$ à $3{,}2~\%$ quel que soit le maillage ; la loi de correction mesurée fixe la pénalité à au moins $0{,}62/\varepsilon$ pour un biais $\varepsilon$. Le contact par pénalité nœud-face, réglage par défaut, fait basculer le bloc glissant au lieu de le laisser glisser ; le contact par potentiel est donc requis dès que deux corps glissent l'un sur l'autre.

La reproduction d'articles est partielle. Sur l'impact d'un insert sur le calcaire de St Anne (Yang et al. 2025~\cite{yang2025validation}), un calcul complet de $300~\mu$s place la contrainte au taillant, la vitesse d'indentation, le rayon du cratère et la fin de la phase de charge à 6 à $12~\%$ des valeurs publiées, avec un bilan d'énergie fermé à $1{,}5\times10^{-7}~\%$. L'enfoncement est trop grand de $17~\%$, la masse de fragments de $87~\%$, et le rebond n'est pas mesuré. Cette comparaison n'est pas indépendante du calage : la loi de joint et la pénalité ont été choisies pendant la même campagne. L'impact sur le granite de Kuru (Yang et al. 2026~\cite{yang2026pulverisation}) n'est pas reproduit : la pulvérisation reste quasi inactive et la loi de joint du code de référence, transcrite telle quelle, crée de l'énergie.

Rien n'est validé au sens strict, c'est-à-dire contre un essai qui n'a pas servi au calage. Les parties III et IV de la campagne de vérification et validation sont vides, sept bancs sont préparés avec leurs critères écrits avant calcul mais n'ont tourné qu'en essai de fumée, et la calibration sur le granite de Red Bohus n'est pas conclue. Les résultats dépendent du nombre de fils de calcul, à environ $2~\%$ sur les pics.
\clearpage
